\SourceLeaf{27}
\textsc{Gauss}' eigene Auseinandersetzungen beziehen sich auf eine specielle Eintheilung der
$(p-1)$ Reste mod.\ $p$ in drei Classen $S$, $S'$, $S''$, wobei die Darstellung von $p$ in der
Form $(mm+3nn)$ benutzt wird. Vermuthlich haben aber die mit dieser Bestimmung der Zahl $f$
verbundenen Umständlichkeiten \textsc{Gauss} davon abgehalten, auf dem hier betretenen Wege
weiterzugehen. In der That wird an dieser Stelle die Nothwendigkeit der Erweiterung des
Zahlgebietes auf die complexen ganzen Zahlen $(a+b\rho)$, wo alsdann an Stelle von $f$ die
Einheitswurzel
$\rho = \frac{-1+i\sqrt{3}}{2}$
selbst tritt, ganz besonders überzeugend.

Die fraglichen Untersuchungen sind jedenfalls in der Zeit vor 1809 ausgeführt und haben sich
wahrscheinlich unmittelbar an die genannten Entwicklungen über quadratische Reste angeschlossen,
welche im Januar 1808 der Gött. Gesellsch. d. Wiss. vorgelegt sind. Nach Notizen in einem von
\textsc{Gauss} seit Ende März 1796 während einer längeren Reihe von Jahren über seine
Untersuchungen geführten Tagebuche\HistoricalNote{\textsuperscript{*)}}{*)}{Über dieses höchst werthvolle Tagebuch sollen in
Bd.\ 9 der ges. Werke nähere Angaben gemacht werden.} ist derselbe Mitte Februar 1807
ausführlicher an die Theorie der cubischen und biquadratischen Reste herangegangen. Andererseits
aber hat sich auch der Ersatz von $f$ durch $\rho$ wahrscheinlich unmittelbar daran
angeschlossen, wie aus der Urschrift der hier vorstehend mitgetheilten Notizen hervorgeht. Man
darf demnach annehmen, dass \textsc{Gauss} bereits um die bezeichnete Zeit (1808) den so
folgenreichen Schritt der Einführung der ganzen complexen Zahlen definitiv vollzogen hat.

\textsc{Gauss} hat auch für die complexen Primzahlen $2-\rho$ und $3-\rho$ die drei Classen $S$,
$S'$, $S''$ aufgestellt (die oben reproducirt sind) sowie gleichfalls noch für den Primtheiler
$3-2\rho$ von 19. Besonders interessant ist, dass \textsc{Gauss} diese Ueberlegungen (wie
entsprechend auch bei den biquadratischen Resten) in geometrische Gestalt kleidete. Die ganzen
complexen Zahlen $a+b\rho$ liefern Punkte der Ebene, welche ein rhombisches Punktgitter
darstellen. Die Gitterpunkte eines einzelnen Systems $S$ lassen sich dabei in einen Bereich
eingrenzen, welcher (im neueren Sinne gesprochen) als Discontinuitätsbereich einer unendlichen
Gruppe linearer Substitutionen aufgefasst werden kann. \textsc{Gauss} hat für verschiedene Fälle
Zeichnungen angefertigt; die für die Primzahl $6-\rho$ ist oben mitgetheilt. Es würde sich hier
um die Gruppe handeln:
\[
u' = \rho^{\nu} u + (a+b\rho),
\]
wo $u$ eine complexe Variabele bedeutet, $\nu = 1, 2, 3$ zu setzen ist und $a+b\rho$ alle der
Congruenz
\[
a+b\rho \equiv 0 \quad(\mathrm{mod.}\ 6-\rho)
\]
genügende ganze Zahlen zu durchlaufen hat. Die in der \textsc{Gauss}'schen Figur angegebenen
Bereiche beziehen sich auf die durch Zusatz der Substitution $u' = \pm u$ erweiterte Gruppe, Die
Willkür
in der Auswahl des Repräsentantensystems $S$ läuft jetzt auf die Willkür in der Auswahl der
Bereichgrenzen hinaus. Welche Absicht \textsc{Gauss} mit der von ihm gewählten Gestalt dieser
Grenzen befolgt hat, ist leider nicht angegeben.

Die am Schlusse der obigen Notizen unter [2] zusammengestellten zwölf Thesen folgen in der
Urschrift direct auf die Betrachtungen der Zahlen $a+b\rho$. Das System aller ganzen complexen
Zahlen $a+b\rho$ erscheint hier eindeutig ersetzt durch das mod.\ $(1+x+xx)$ reducirte System
aller ganzen Functionen von $x$ mit ganzen rationalen Coefficienten.
\par\nobreak\noindent\hfill\textsc{Fricke}.\par
\endgroup
